\section{Audit, integration status, and further research}
\label{bmh:sec:audit}

\subsection{What is proved and what is unchanged}
The principal analytic results of this article are exact ordinary
integral identities. Their proofs establish the necessary convergence,
parameter differentiation, cancellation, and finite coefficient
extraction. The delivered finite tests check implementations of these
arguments; they are not a substitute for the arguments.

The sampled source material did not reveal a new false equation in the
one-factor theorem or the cited canonical depth results. Accordingly we
do not manufacture a mathematical erratum. The concrete repository
change is a \emph{status refinement} for the tenth research question in
\src{sections/06-audit-research.tex}: its positive-integer-order finite-
closure sector now has a constructive solution, while the arbitrary-
spectral-order and minimal-depth portions remain open.

The current Gaussian $S_6$ and $S_8$ identities are not proved or
refuted by this work. The prior exact $S_4$ result is not counted again
as a new result. No claim is made that the binary incoming archives have
all been read, that all new claims in those archives have been reconciled,
or that these proofs have been checked in a proof assistant.

\subsection{Normalization corrections to avoid during integration}
The following are explicit safeguards for the new continuation, not
accusations that the inspected source made each mistake.
\begin{enumerate}[leftmargin=2em,itemsep=5pt]
\item \textbf{Use the product strip.} For two factors the lower condition
is $\Rea a>-2$, not $\Rea a>-1$. The value at $a=-1$ is an ordinary
integral and has the exact formula~\eqref{bmh:eq:negative-jets}.
\item \textbf{Use connected leading-one coordinates.} The separate
objects $\zeta(1,B)$ and $Z_q(1,B)$ generally diverge. Equations
\eqref{bmh:eq:Cedge}--\eqref{bmh:eq:Ccorner} give convergent replacements,
derived by a combined analytic continuation.
\item \textbf{Keep the zero-order compact value correct.}
$\Li_0(-t/(1-t))=-t$. Using $-t/(1-t)$ breaks the kernel recurrence,
including its simplest beta base.
\item \textbf{Cancel poles before differentiating a specialization.}
Examples include $a=0,-1$, $s=1$, and the one-factor product
$s\zeta(s+1)$ at $s=0$. The last removable value is $1$, not $0$.
\item \textbf{Separate algebra membership from additive depth.}
Even-zeta factors times double-zeta values belong to the algebra generated
by doubles, but can have additive depth three. The word-weight filtration
likewise supplies no theorem of numerical independence.
\item \textbf{Do not call a discrete primitive an antiderivative.}
Equation~\eqref{bmh:eq:Stieltjes-finite-sum} is a shift identity;
\eqref{bmh:eq:continuous-primitive} and~\eqref{bmh:eq:primitive} are
continuous anchored primitives with explicit constants.
\end{enumerate}

\subsection{Proposed research questions}
The questions below seek exact identities or structural classifications,
not numerical lower bounds inferred from failed searches.
\begin{enumerate}[label=\textbf{Q\arabic*.},leftmargin=2.8em,itemsep=7pt]
\item \textbf{A genuinely two-scale double-Hurwitz formula.}
For independent $z,w$, replace the unit-scale cone decomposition by a
shifted-cone calculation and identify a minimal convenient family of
multiple Lerch coordinates. The answer should specialize to
\eqref{bmh:eq:master} when $z=w=1$ and to
\eqref{bmh:eq:multiscale-example} at logarithmic order.
\item \textbf{Three or more equal-scale factors.}
A product of three Fermi integrals splits into six order cones. Determine
a finite connected triple-Hurwitz system that makes all leading-one
faces convergent. The two-factor coordinates suggest an inclusion--
exclusion pattern, but its exact signs and edge terms require proof.
\item \textbf{Depth of the weighted MZV family.}
The specific weighted sums in~\eqref{bmh:eq:weighted} reduce to double
values at every weight. Find a short double-shuffle proof that does not
pass through the integral, and determine which further coefficient
symmetries force a single-zeta reduction. Any arithmetic minimal-depth
claim requires additional input and must be stated separately.
\item \textbf{The missing central parity sector.}
When $W+h$ is even, inversion removes the desired central moment from
\eqref{bmh:eq:reciprocity}. Determine an independent functional relation
for this complementary sector. Further differentiation of the same
projection cannot recover a coefficient that it annihilates.
\item \textbf{A complete continuous primitive system.}
Equation~\eqref{bmh:eq:continuous-primitive} integrates a specified
adjacent-index combination. Construct a basis in which individual
connected coordinates have explicit primitives, with all integration
constants and leading-one specializations controlled.
\item \textbf{Spectral derivatives of the double-Hurwitz master.}
The present cone expansion uses integer $m,n$. Find a representation at
complex $m,n$ that is differentiable in the spectral variables and
specializes to the finite binomial formula. This is a precise route toward
products containing $\Li_s^{[r]}$ without assuming finite ordinary-MPL
closure at noninteger order.
\item \textbf{Rational common shifts and cyclotomic coordinates.}
At $q=p/d>0$, express the connected double-Hurwitz values and their
Stieltjes spectral jets through residue-class filters. Quantify the finite
alphabet and the exact cancellation at leading-one coordinates before
reducing Dirichlet-character combinations.
\item \textbf{Higher negative resonances for more factors.}
An $r$-factor product has ordinary integer values at
$a=1-r,\ldots,-1$. Determine whether the connected higher-depth system
has a shift law as simple as~\eqref{bmh:eq:Cshift}, and derive every
logarithmic moment at those resonances.
\item \textbf{A proof-assistant boundary.}
The closed shuffle multiplicities, Laurent decomposition, finite
recurrences, and coefficient convolutions are natural first targets for
formalization. The remaining analytic obligations are explicit: dominated
holomorphy, the Fermi integral, the combined cone limit, and endpoint
regularization. No finite test count alone discharges these obligations.
\end{enumerate}

\subsection{Integration recommendation}
Place this delivery as a continuation of the Mellin--Lerch report under
\src{continuations/bilinear-mellin-hurwitz/}, or under a report path chosen
by the repository intake process. Preserve the self-contained article and
its checks. The canonical integration chapter can adopt the master,
resonance, and kernel sections after analytic review; the differentiation
chapter can adopt the Stieltjes identities, and the depth chapter can
adopt the weighted MZV theorem. Every theorem and equation label here has
the prefix \src{bmh:}. The accompanying integration note supplies the
precise scope wording for the research-question update.
